\begin{proof}
\begin{enumerate}[label=\arabic*.]
\item \textbf{The actual height, parameter, and real coordinate maps.}
Put $k=625\pi/96$, $c_0=k^{-1}$, and use the specified right
connection $H'=H\Gamma_6$, $H(1)=I$, with $Y=XH$.
For the uniquely normalized real quotient-intersection vector write
\[
 q=(0,v_2,v_1,-i,0,0)^t,\qquad z=Yq,
 \qquad h_b=c_0\operatorname{Im}z_0.
\]
The lower five entries of $z$ are real. Every marked monodromy
matrix $T$ is real and fixes $e_0$. Its top entry on $z$ is
$z_0$ plus a real linear combination of those lower entries, so
$\operatorname{Im}(Tz)_0=\operatorname{Im}z_0$. This proves
branch independence in the given marking.
The normalization gives $\alpha(q)=k$, since the coefficient of
$w_0$ is $-i$ and $\alpha(w_0)=ki$. The lifted vector for parameter
$\xi$ therefore has top coordinate $z_0+k\xi$ and the same lower
coordinates. The exact obstruction in the real quotient-intersection
line is
\[
 \operatorname{Im}(z_0+k\xi)=k(h_b+\operatorname{Im}\xi).
\]
The transversality criterion of \cref{night:P1.8} consequently
uses this same function for complex and real parameters.

The numerical preconditioner has an explicit real coordinate map.
For a constant limiting root frame $H_c$, let $B$ have lower-five
columns
\[
 \operatorname{Im}(H_c)_1,\quad\operatorname{Re}(H_c)_1,
 \quad\operatorname{Im}(H_c)_2,\quad\operatorname{Re}(H_c)_2,
 \quad\operatorname{Im}(H_c)_3.
\]
The certified inverse of $B$ defines
\[
 r=(\operatorname{Im}(H_c)_{01},\operatorname{Re}(H_c)_{01},
       \operatorname{Im}(H_c)_{02},\operatorname{Re}(H_c)_{02},0)B^{-1},
 \qquad
 T_c=\begin{pmatrix}1&-r\\0&B^{-1}\end{pmatrix}.
\]
This is a real invertible matrix fixing $e_0$.
Realness of $B^{-1}z_{\rm low}$ is equivalent to realness of
$z_{\rm low}$, and the change of top entry is real on that vector.
Thus it preserves both the normalized pair and $h_b$ exactly.
Writing $v_2=a+ib$, $v_1=c+id$ gives the actual real equations
with row
$(\operatorname{Im}Y_{i1},\operatorname{Re}Y_{i1},
\operatorname{Im}Y_{i2},\operatorname{Re}Y_{i2})$
and right-hand entry $\operatorname{Re}Y_{i3}$.
Approximate row selection only chooses four rows; the inverse
and every ensuing enclosure are certified in ball arithmetic.

\item \textbf{The marked conjugation map.}
The connection has real rational coefficients. Conjugating the
specified ordered polygonal meridians gives their inverse paths.
The actual matrix $K_0=X\overline X^{-1}$ therefore satisfies
$K_0A=A^{-1}K_0$, $K_0B=B^{-1}K_0$, and $K_0e_0=e_0$.
The normalized intertwiner uniqueness in \cref{night:P1.15}
applies: the quotient of two invertible solutions is in the
common centralizer fixing $e_0$. Direct exact solution gives
\[
 K=\begin{pmatrix}
 1&2&-1&-1&2&-1\\
 0&0&0&0&-1&1\\
 0&0&-1&0&0&1\\
 0&0&0&-1&0&1\\
 0&-1&0&0&0&1\\
 0&0&0&0&0&1
 \end{pmatrix}.
\]
Multiplication verifies the two intertwiner equations,
$Ke_0=e_0$, $K^2=I$, and $\det K=-1$, hence $K=K_0$.
On conjugate continuation branches,
$Y(\overline\lambda)=K\overline{Y(\lambda)}$.
The vector $q'=-\overline q$ has the required $-i$ normalization,
and its image is $-K\overline z$, with real lower coordinates.
Uniqueness therefore identifies the normalized intersection
vector and gives
$h_b(\overline\lambda)=h_b(\lambda)$.
This is the actual coordinate map used in reflecting the covered
regions.

\item \textbf{Rigorous analytic enclosures.}
The following formulas specify the complete-domain computation.
All matrix coefficients come from the displayed rational
connection and the specified gauges of
\cref{night:P1.15,night:P1.16}. A regular root frame is the
pullback followed by its stated diagonal gauge $D$; its right
connection is exactly
$D^{-1}\lambda'\Gamma_6D+D^{-1}D'$.
At a positive anchor the constant limiting frame is obtained by
the actual map $H_c=Y_{\rm root}(z_c)U(z_c)^{-1}$, using the
normalized continuation and a certified matrix inverse.
